\section{Average curvature and index rigidity}\label{sec:index}

Throughout this section, $\Sigma^m\subset S^{m+1}$ is a smooth, closed,
connected, embedded minimal hypersurface, $m\ge2$, which is not totally
geodesic.  Write $x$ for its position vector, $N$ for a global unit normal,
$A=-dN$ for its shape operator, and $P=|A|^2$.  Integrals use the induced
volume measure $dV$; set $V_\Sigma=\Vol_m(\Sigma)$.  The inequality
\[
  \int_\Sigma P\,dV\ge mV_\Sigma
\]
is the inequality part of Perdomo's conjecture, as formulated in
\cite[Conjecture~1.5]{GeLi2021}.  We prove it by a two-sided curvature
maximum.  Combining it with the conformal and spectral argument of
\cite[Theorem~3.1]{Perdomo2004} then gives the index rigidity needed below.
We include that argument, in particular its pointwise equality step.

With $\Delta=\operatorname{div}\nabla$, the minimal hypersurface identities
are
\begin{equation}\label{eq:index-geometric-identities}
 \Delta x=-mx,\qquad \Delta N=-PN,\qquad
 \Delta A=(m-P)A,\qquad
 \frac12\Delta P=|\nabla A|^2+(m-P)P.
\end{equation}
Here the Laplacian of $A$ is the rough tensor Laplacian.  We also use
$\tr A=0$ and the Codazzi identity: $\nabla A$, viewed as a covariant
three-tensor, is symmetric.  The last two formulas are the hypersurface
case of Simons' identity
\cite[Theorem~5.3.1, equation~(5.3.2)]{Simons1968}.

\subsection{A two-sided curvature maximum}

Define
\begin{equation}\label{eq:index-mu}
 \mu(x)=\sup_{\substack{y\in\Sigma\\y\ne x}}
 \frac{|\langle N(x),y\rangle|}{1-\langle x,y\rangle}.
\end{equation}
Embeddedness makes the denominator positive away from the diagonal.
The absolute value controls both signs of the principal curvatures.
The reflection used below appears in Brendle's two-point argument
\cite[Section~2]{Brendle2013} and in curvature estimates for hypersurface
flows \cite[Section~2]{AHLW2015}. The diagonal and distinct-point contact
calculations are also closely related to
\cite[Lemma~3.5 and Proposition~4.1]{Nifa2026}.
Using the sum of the two one-sided tangent-ball curvature functions,
Nifa obtains a stronger quadratic-form inequality and derives both the
average-curvature bound and the embedded least-index classification
\cite[Theorem~1.1 and Corollary~7.3]{Nifa2026}.
For the integrated estimate needed here, the absolute maximum in
\eqref{eq:index-mu} suffices. We prove its differential inequality directly.

Rather than differentiate the maximum, we estimate smooth upper tests.
The logarithmic form of the estimate will allow comparison with a Poisson
potential, forcing the mean of $m-P$ to be nonpositive.

\begin{lemma}[Upper tests for the curvature maximum]
\label{lem:index-upper-tests}
The function $\mu$ in \eqref{eq:index-mu} is finite, continuous, and strictly
positive.  If a smooth positive function $\phi$ satisfies $\phi\ge\mu$ in
a neighborhood of $x$ and $\phi(x)=\mu(x)$, then
\begin{equation}\label{eq:index-log-test}
 \Delta\log\phi(x)\ge m-P(x).
\end{equation}
\end{lemma}

\begin{proof}
We first describe the possible contacts.  Let
$\gamma(t)=\exp_x(tv)$, where $v\in T_x\Sigma$ is a unit vector, and regard
$\gamma$ as a curve in $\mathbb R^{m+2}$.  The ambient differentiation
formulas give
\begin{align}
 \langle N(x),\gamma(t)\rangle
   &=\frac{t^2}{2}A_x(v,v)
     +\frac{t^3}{6}(\nabla_vA)_x(v,v)+O(t^4),
       \label{eq:index-contact-numerator}\\
 1-\langle x,\gamma(t)\rangle
   &=\frac{t^2}{2}+O(t^4).
       \label{eq:index-contact-denominator}
\end{align}
For example, $\gamma''(0)=-x+A_x(v,v)N(x)$, while
$\langle x,\gamma'''(0)\rangle=0$.  Compactness makes these expansions
uniform on the unit tangent bundle for sufficiently small $|t|$.
Consequently, adjoining unit tangent directions at the diagonal extends
the quotient continuously, with values $|A_x(v,v)|$.

For completeness, this extension also gives continuity when the base
point varies.  If $x_j\to x$ and nearly maximizing points $y_j$ stay away
from $x$, a convergent subsequence gives an ordinary distinct-point
contact.  If $y_j\to x$, write
$y_j=\exp_{x_j}(t_jv_j)$ and take a convergent subsequence of unit tangent
directions.  The uniform expansions show that its limiting quotient is
at most $\mu(x)$.  This proves upper semicontinuity.  Holding fixed any
$y\ne x$ whose quotient is within $\varepsilon$ of $\mu(x)$ proves lower
semicontinuity, after letting $\varepsilon\downarrow0$.  The same
compactness argument proves finiteness and shows that a maximum is
realized either by distinct points or by a tangent direction.  If
$\mu(x)=0$, the entire hypersurface lies in $N(x)^\perp$, contrary to
fullness from Lemma~\ref{lem:fullness}.  Thus $\mu>0$.

Fix an upper test $\phi$ at $x$.  At a realizing contact, choose the sign
of $N$ to make the signed quotient positive, and change the sign of $A$
at the same time.  This choice is held fixed in a neighborhood.  It
requires no continuous choice of maximizing sign: for either fixed sign
and for every nearby $x'$ and every $y'\ne x'$,
\begin{equation}\label{eq:index-fixed-sign}
 \frac{\langle N(x'),y'\rangle}{1-\langle x',y'\rangle}
 \le\mu(x')\le\phi(x').
\end{equation}
Choose a principal orthonormal basis $e_1,\ldots,e_m$ at $x$, and write
$Ae_i=\lambda_i e_i$.  The diagonal limits give
\begin{equation}\label{eq:index-two-sided-eigenvalues}
 -\phi(x)\le\lambda_i\le\phi(x)\qquad(1\le i\le m).
\end{equation}
We shall prove the stronger inequality
\begin{equation}\label{eq:index-gradient-test}
 \Delta\phi+(P-m)\phi
 \ge 2\sum_{\lambda_i<\phi}
             \frac{\phi_i^2}{\phi-\lambda_i}
 \ge\frac{|\nabla\phi|^2}{\phi}
 \quad\text{at }x,
\end{equation}
where $\phi_i=d\phi(e_i)$ and every component with $\lambda_i=\phi$
will be shown to vanish.

\smallskip
\noindent\emph{Distinct-point contact.}
Suppose $y\ne x$ realizes the maximum with the chosen sign.  On the
product of a neighborhood of $x$ with $\Sigma$, set
\[
 Z(x',y')=\phi(x')\bigl(1-\langle x',y'\rangle\bigr)
                  -\langle N(x'),y'\rangle.
\]
By \eqref{eq:index-fixed-sign}, $Z\ge0$ and $Z(x,y)=0$.  Write
$c=\langle x,y\rangle$, $s=1-c>0$, and $d_i=\langle e_i,y\rangle$.
The first derivatives at the minimum give
\begin{equation}\label{eq:index-distinct-first}
 s\phi_i=(\phi-\lambda_i)d_i.
\end{equation}
Reflection across the hyperplane perpendicular to $x-y$ sends $x$ to
$y$, and sends $N(x)$ to
\[
 n_y=N(x)+\phi(x-y).
\]
Indeed $\langle N(x),y\rangle=\phi s$, so this vector is perpendicular to
$y$ and has length one.  The $y$ first-derivative equation for $Z$ shows
that $n_y$ is perpendicular to $T_y\Sigma$.  Therefore the reflected
vectors
\[
 \widetilde e_i=e_i+\frac{d_i}{s}(x-y)
\]
form an orthonormal basis of $T_y\Sigma$, and
\[
 \langle x,\widetilde e_i\rangle=d_i,
 \qquad
 \langle e_i,\widetilde e_i\rangle=1-\frac{d_i^2}{s}.
\]
The vector $n_y$ may be either orientation of the normal at $y$; only
minimality at $y$ enters the following trace computation.

Use the product connection to compute the Hessian of $Z$.  Denote its
mixed entries in the indicated bases by $Z_{x_i y_i}$.  From
\eqref{eq:index-geometric-identities},
\begin{align*}
 \Delta_x Z
   &=s\Delta\phi-2\sum_i\phi_i d_i+m\phi c+P\phi s,\\
 \Delta_y Z&=m\phi,\\
 \sum_i Z_{x_i y_i}
   &=-\sum_i\phi_i d_i
     -\sum_i(\phi-\lambda_i)
                         \left(1-\frac{d_i^2}{s}\right).
\end{align*}
In particular, the nonnegative sum of Hessians along
$(e_i,\widetilde e_i)$ is
\begin{align*}
 0&\le\Delta_xZ+2\sum_iZ_{x_i y_i}+\Delta_yZ\\
  &=s\bigl(\Delta\phi+(P-m)\phi\bigr)
       -4\sum_i\phi_i d_i
       +\frac{2}{s}\sum_i(\phi-\lambda_i)d_i^2\\
  &=s\bigl(\Delta\phi+(P-m)\phi\bigr)
       -\frac{2}{s}\sum_i(\phi-\lambda_i)d_i^2,
\end{align*}
where we used $\tr A=0$ and then
\eqref{eq:index-distinct-first}.  If $\lambda_i=\phi$,
\eqref{eq:index-distinct-first} gives $\phi_i=0$.  Dividing by $s$ and
substituting that equation in all other directions gives the first
inequality of \eqref{eq:index-gradient-test}.

\smallskip
\noindent\emph{Diagonal contact, including repeated eigenvalues.}
Suppose instead that a unit tangent direction realizes the maximum.
With our choice of normal sign, the top eigenvalue is $\phi(x)$.
Let $E\subset T_x\Sigma$ be its full eigenspace.  For every unit $v\in E$
and every local smooth unit extension $V$ of $v$,
\[
 A(V,V)\le\mu\le\phi,
 \qquad A_x(v,v)=\phi(x).
\]
The first derivative at this contact, in any direction $w\in T_x\Sigma$,
is
\begin{equation}\label{eq:index-diagonal-first}
 d\phi(w)=(\nabla_wA)(v,v).
\end{equation}
The field-derivative term vanishes because
$A(v,\nabla_wV)=\phi\langle v,\nabla_wV\rangle=0$.
For the geodesic with initial direction $v$, formulas
\eqref{eq:index-contact-numerator}--\eqref{eq:index-contact-denominator}
give the signed quotient
\[
 \frac{\langle N(x),\exp_x(tv)\rangle}
      {1-\langle x,\exp_x(tv)\rangle}
 =\phi(x)+\frac{t}{3}(\nabla_vA)(v,v)+O(t^2).
\]
It is at most the fixed value $\mu(x)=\phi(x)$ for both signs of $t$.
Hence $(\nabla_vA)(v,v)=0$, and
\eqref{eq:index-diagonal-first} implies
\begin{equation}\label{eq:index-top-gradient}
 d\phi|_E=0.
\end{equation}
This holds for every vector in $E$, not just for one principal vector.
Moreover, polarizing \eqref{eq:index-diagonal-first} on $E$ gives
$(\nabla_wA)|_{E\times E}=d\phi(w)g|_{E\times E}$.
These are pointwise statements; no smooth extension of the eigenspace
is assumed.

Fix $e_1\in E$ and choose a smooth unit field $V$ with $V(x)=e_1$.
Put
\[
 T_{ji}=(\nabla_{e_j}A)(e_1,e_i),\qquad
 \beta_{ji}=\langle\nabla_{e_j}V,e_i\rangle.
\]
The first derivatives $\nabla_{e_j}V$ may be prescribed arbitrarily in
$e_1^\perp$: prescribe the corresponding jets of any local field and
normalize it.  Unit length determines precisely the second-derivative
contraction that occurs in the calculation,
\[
 \langle e_1,\Delta V\rangle
       =-\sum_j|\nabla_{e_j}V|^2.
\]
Thus, in a geodesic orthonormal frame at $x$,
\begin{equation}\label{eq:index-unit-field-laplacian}
 \Delta\bigl(A(V,V)\bigr)
  =(\Delta A)_{11}
    +4\sum_{j,i}\beta_{ji}T_{ji}
    +2\sum_{j,i}(\lambda_i-\phi)\beta_{ji}^2.
\end{equation}
No other second derivative of $V$ contributes, since $Ae_1=\phi e_1$.
Choose the free first jets by
\[
 \beta_{ji}=\frac{T_{ji}}{\phi-\lambda_i}
       \quad\text{if }\lambda_i<\phi,
 \qquad
 \beta_{ji}=0\quad\text{if }\lambda_i=\phi.
\]
The minimum of $\phi-A(V,V)$ then yields
\begin{align*}
 \Delta\phi
 &\ge(\Delta A)_{11}
       +2\sum_{\lambda_i<\phi}
                    \frac{\sum_jT_{ji}^2}{\phi-\lambda_i}\\
 &\ge(m-P)\phi
       +2\sum_{\lambda_i<\phi}
                    \frac{\phi_i^2}{\phi-\lambda_i}.
\end{align*}
In the last step we used Simons' identity and Codazzi:
$T_{1i}=(\nabla_{e_i}A)(e_1,e_1)=\phi_i$.
Equation~\eqref{eq:index-top-gradient} accounts for every omitted
top-eigenspace gradient component.  This proves the first inequality of
\eqref{eq:index-gradient-test} also at a repeated diagonal contact.

In both cases, \eqref{eq:index-two-sided-eigenvalues} gives
$0<\phi-\lambda_i\le2\phi$ in each retained direction.  The vanished
components in the other directions therefore give the second inequality
of \eqref{eq:index-gradient-test}.  Dividing by $\phi$ and subtracting
$|\nabla\phi|^2/\phi^2$ proves \eqref{eq:index-log-test}.
\end{proof}

\begin{theorem}[Average-curvature inequality]\label{thm:curvature-average}
Every smooth, closed, connected, embedded minimal hypersurface
$\Sigma^m\subset S^{m+1}$, $m\ge2$, which is not totally geodesic satisfies
\begin{equation}\label{eq:index-curvature-average}
 \int_\Sigma |A|^2\,dV\ge m\Vol_m(\Sigma).
\end{equation}
\end{theorem}

\begin{proof}
If the inequality failed, then
\[
 c=\frac{1}{V_\Sigma}\int_\Sigma(m-P)\,dV>0.
\]
Since $m-P-c$ has mean zero, the Poisson equation on the closed connected
manifold has a smooth solution. Indeed, the kernel of $\Delta$ consists
of constants by integration by parts; the compact elliptic Fredholm
alternative and regularity
\cite[\S10, equations~(10.4)--(10.6)]{TaylorEllipticNotes} give
\begin{equation}\label{eq:index-poisson}
 \Delta h=m-P-c.
\end{equation}
Let $K=\max_\Sigma(\log\mu-h)$ and let $x$ attain this maximum.
The smooth positive function $\phi=\exp(h+K)$ is an upper test for $\mu$
at $x$.  Lemma~\ref{lem:index-upper-tests} and
\eqref{eq:index-poisson} imply
\[
 m-P(x)\le\Delta\log\phi(x)=m-P(x)-c,
\]
a contradiction.  This argument uses only continuity and positivity of
$\mu$; it does not differentiate $\mu$ or assume any weak regularity for
its gradient.
\end{proof}

\subsection{Normal coordinates and conformal tests}

The average-curvature bound is now available. To turn it into index
rigidity, we need conformal coordinates orthogonal to the top Jacobi
eigenfunction and a strict bound for their energy. Equality of the
curvature integrals alone will not be used as a pointwise curvature
statement: that step will follow from equality in each spectral test.

\begin{lemma}[Independence of the normal coordinates]
\label{lem:index-normal-independence}
The map
$a\mapsto\langle a,N\rangle$ from $\mathbb R^{m+2}$ to
$C^\infty(\Sigma)$ is injective.
\end{lemma}

\begin{proof}
Suppose $a\ne0$ and $\langle a,N\rangle=0$ everywhere.  Set
$f=\langle a,x\rangle$.  The ambient Gauss formula gives
\[
 \nabla^2f=-fg+\langle a,N\rangle A=-fg.
\]
This function is nonconstant: a constant solution would be zero, whereas
$f\equiv0$ contradicts fullness.  Obata's sphere-rigidity theorem
\cite[Theorem~A]{Obata1962} applies to the complete connected manifold $(\Sigma,g)$
and shows that it is isometric to the unit round $S^m$.  The scalar Gauss
equation for a minimal hypersurface,
$\operatorname{Scal}_\Sigma=m(m-1)-P$, then gives $P=0$.  This contradicts
the hypothesis that $\Sigma$ is not totally geodesic.
\end{proof}

For $a\in S^{m+1}$ and $\tau\ge0$, put $C=\cosh\tau$, $S=\sinh\tau$,
and define a conformal automorphism of the ambient sphere by
\begin{equation}\label{eq:index-conformal-map}
 T_{\tau,a}(X)=
 \frac{X+\bigl((C-1)\langle a,X\rangle-S\bigr)a}
      {C-S\langle a,X\rangle}.
\end{equation}
Its length factor is
\begin{equation}\label{eq:index-conformal-factor}
 \beta_{\tau,a}(X)=\bigl(C-S\langle a,X\rangle\bigr)^{-1}.
\end{equation}
Indeed, with $B=I+(C-1)aa^{\mathsf T}$, differentiation gives
\[
 dT_{\tau,a}(W)=\beta_{\tau,a}
       \bigl(BW+S\langle a,W\rangle T_{\tau,a}(X)\bigr),
 \qquad W\perp X,
\]
whose inner products are $\beta_{\tau,a}^2$ times those of the original
tangent vectors.  In particular $T_{0,a}$ is the identity.

The balancing step is the conformal center-of-mass argument of Li and
Yau \cite[proof of Theorem~1, p.~274]{LiYau1982}, applied to the positive
Jacobi eigenfunction as in Perdomo's argument
\cite[Theorem~3.1]{Perdomo2004}. We include the construction and the strict
area identity that will be needed for equality.

\begin{lemma}[Conformal balancing and strict area maximum]
\label{lem:index-conformal-tests}
For every smooth positive weight $u$ on $\Sigma$, there is a map
$T_{\tau,a}$ in \eqref{eq:index-conformal-map} such that
\begin{equation}\label{eq:index-balance}
 \int_\Sigma u\,T_{\tau,a}(x)\,dV=0.
\end{equation}
For every $a\in S^{m+1}$ and $\tau\ge0$,
\begin{equation}\label{eq:index-conformal-area}
 \int_\Sigma\beta_{\tau,a}^{\,m}\,dV\le V_\Sigma,
\end{equation}
with strict inequality if $\tau>0$.  Consequently,
\begin{equation}\label{eq:index-conformal-energy}
 \int_\Sigma\beta_{\tau,a}^{\,2}\,dV\le V_\Sigma,
\end{equation}
again with strict inequality if $\tau>0$.
\end{lemma}

\begin{proof}
Parametrize the transformations by $z=(\tanh\tau)a$ in the open unit ball
of $\mathbb R^{m+2}$, with the identity at $z=0$.  Consider
\[
 F(z)=\frac{\int_\Sigma u\,T_{\tau,a}(x)\,dV}
                  {\int_\Sigma u\,dV}.
\]
As $z\to a_*\in S^{m+1}$, the integrand tends to $-a_*$ at every point
with $x\ne a_*$.  This remains true for a moving direction $a\to a_*$.
The measure $u\,dV$ assigns zero mass to a point, and the integrands have
length one.  Dominated convergence extends $F$ continuously to the
closed ball, with $F(a_*)=-a_*$.  If $F$ had no zero, its normalization
$F/|F|$ would extend the antipodal map from the boundary sphere to the
ball.  The antipodal map has nonzero degree, so such an extension is
impossible.  A zero lies in the open ball and proves
\eqref{eq:index-balance}.

To prove \eqref{eq:index-conformal-area}, fix $a$ and abbreviate
$f=\langle a,x\rangle$, $\eta=\langle a,N\rangle$, and
$\beta=\beta_{\tau,a}(x)$.  Orthogonal decomposition of $a$ into position,
tangent, and normal components gives
\[
 |\nabla f|^2=1-f^2-\eta^2,\qquad
 \Delta f=-mf,\qquad \nabla\beta=S\beta^2\nabla f.
\]
Integrating a divergence on the closed hypersurface yields
\begin{align*}
 0&=\frac1m\int_\Sigma\operatorname{div}(\beta^m\nabla f)\,dV\\
  &=\int_\Sigma\beta^{m+1}
            \bigl(S(1-f^2-\eta^2)-f(C-Sf)\bigr)\,dV\\
  &=\int_\Sigma\beta^{m+1}(S-Cf-S\eta^2)\,dV.
\end{align*}
Differentiation in $\tau$ therefore gives the exact formula
\begin{equation}\label{eq:index-area-derivative}
 \frac{d}{d\tau}\int_\Sigma\beta^m\,dV
   =m\int_\Sigma\beta^{m+1}(Cf-S)\,dV
   =-mS\int_\Sigma\beta^{m+1}\eta^2\,dV.
\end{equation}
By Lemma~\ref{lem:index-normal-independence}, $\eta$ is not identically
zero.  The derivative is strictly negative for every $\tau>0$.
Starting from $\beta_{0,a}=1$ proves \eqref{eq:index-conformal-area}
and its strictness.  Finally, H\"older's inequality, including the
identity case $m=2$, gives
\[
 \int_\Sigma\beta^2\,dV
 \le V_\Sigma^{\,1-2/m}
          \left(\int_\Sigma\beta^m\,dV\right)^{2/m}
 \le V_\Sigma.
\]
The second inequality is strict for $\tau>0$, proving
\eqref{eq:index-conformal-energy}.
\end{proof}

\subsection{The least possible index}

We combine the curvature lower bound with the conformal energy upper
bound. Under the least-index hypothesis, the two estimates force
equality in the spectral comparisons and hence parallel second
fundamental form.

For our Laplacian convention, the Jacobi operator and the normal second
variation are
\begin{equation}\label{eq:index-jacobi}
 L=\Delta+m+P,\qquad
 \delta^2\Area(\phi N,\phi N)
   =-\int_\Sigma\phi L\phi\,dV
   =\int_\Sigma\bigl(|\nabla\phi|^2-(m+P)\phi^2\bigr)\,dV.
\end{equation}
Thus $\ind(\Sigma)$ counts the positive eigenvalues of $L$, with
multiplicity. The operator is self-adjoint on $H^2(\Sigma)$ and has a
complete orthonormal basis of smooth eigenfunctions
\cite[\S10, equations~(10.7)--(10.11)]{TaylorEllipticNotes}.
Its spectrum is bounded above, since its quadratic form is at most
$\max_\Sigma(m+P)$ times the squared $L^2$ norm.

\begin{proposition}[Index rigidity]\label{prop:index-rigidity}
Every smooth, closed, connected, embedded minimal hypersurface
$\Sigma^m\subset S^{m+1}$, $m\ge2$, which is not totally geodesic has
\[
 \ind(\Sigma)\ge m+3.
\]
If $\ind(\Sigma)=m+3$, then, up to an ambient isometry, its image is
\[
 S^k\!\left(\sqrt{\frac{k}{m}}\right)
 \times
 S^{m-k}\!\left(\sqrt{\frac{m-k}{m}}\right)
 \quad\text{for some }1\le k<m.
\]
In particular, such a hypersurface of volume less than $a_m$ has
index strictly greater than $m+3$.
\end{proposition}

\begin{proof}
By \eqref{eq:index-geometric-identities}, the $m+2$ normal coordinate
functions satisfy $LN=mN$; they are linearly independent by
Lemma~\ref{lem:index-normal-independence}.  The largest eigenvalue
$\lambda_1$ of $L$ is simple and has a strictly positive eigenfunction
$u$. Apply the first-eigenfunction theorem to the shifted Schr\"odinger
operator $K-L$, with $K>1+\max_\Sigma(m+P)$; its positivity uses
connectedness and the strong maximum principle
\cite[Proposition~2.9]{TaylorSpectralNotes}. Testing its Rayleigh quotient with
the constant function gives
\[
 \lambda_1\ge m+\frac{1}{V_\Sigma}\int_\Sigma P\,dV>m.
\]
These eigenfunctions prove $\ind(\Sigma)\ge m+3$.

Suppose that equality holds.  The positive spectrum consists of the
simple eigenvalue $\lambda_1>m$ and the eigenvalue $m$ with multiplicity
$m+2$.  The spectral theorem therefore gives
\begin{equation}\label{eq:index-spectral-bound}
 \int_\Sigma\phi L\phi\,dV\le m\int_\Sigma\phi^2\,dV
 \qquad\text{whenever }\int_\Sigma u\phi\,dV=0.
\end{equation}
Choose the conformal map supplied by
Lemma~\ref{lem:index-conformal-tests} for the weight $u$, and write
$v=T_{\tau,a}(x)$.  Each coordinate $v_\alpha$ satisfies the
orthogonality in \eqref{eq:index-spectral-bound}.  Since $|v|=1$ and
$\sum_\alpha|\nabla v_\alpha|^2=m\beta^2$, summing those inequalities
and integrating by parts gives
\begin{equation}\label{eq:index-spectral-sum}
 \int_\Sigma P\,dV
       \le m\int_\Sigma\beta^2\,dV\le mV_\Sigma.
\end{equation}
Theorem~\ref{thm:curvature-average} forces equality throughout.  The
strict part of \eqref{eq:index-conformal-energy} implies $\tau=0$, so
$v=x$.

We now use equality in the individual spectral inequalities.  For each
$\alpha$, the defect
\[
 d_\alpha=m\int_\Sigma x_\alpha^2\,dV
                  -\int_\Sigma x_\alpha Lx_\alpha\,dV
\]
is nonnegative by \eqref{eq:index-spectral-bound}, and
\eqref{eq:index-spectral-sum} shows $\sum_\alpha d_\alpha=0$.
Thus every defect vanishes.  In the spectral decomposition on
$u^\perp$, all eigenvalues are at most $m$; equality in its quadratic
form implies $Lx_\alpha=mx_\alpha$ for every $\alpha$.  On the other hand,
\eqref{eq:index-geometric-identities} gives $Lx=Px$.  Taking the inner
product of $(P-m)x=0$ with $x$ yields
\begin{equation}\label{eq:index-pointwise-equality}
 P=m\quad\text{everywhere on }\Sigma.
\end{equation}
The pointwise conclusion uses spectral equality as well as the
integrated curvature equality.

The last identity in \eqref{eq:index-geometric-identities} now gives
$\nabla A=0$.  We finish by deriving the product classification.
Parallelness makes the eigenvalues constant and their eigendistributions
parallel.  If $U$ and $W$ are unit vectors in distinct eigendistributions,
the curvature tensor preserves each of those distributions, so
$\langle R(U,W)W,U\rangle=0$.  The Gauss equation gives
\[
 1+\lambda\kappa=0
\]
for any two distinct eigenvalues $\lambda,\kappa$.  There are at most
two distinct eigenvalues: a third would have to equal $-1/\lambda$ as
well.  Since $A$ is nonzero and trace-free, there are exactly two, of
opposite signs.  After changing $N$ if necessary, write them as
\[
 \lambda=\sqrt{\frac{\ell}{k}},\qquad
 \kappa=-\frac1\lambda=-\sqrt{\frac{k}{\ell}},
 \qquad k+\ell=m,\quad k,\ell>0,
\]
where $k$ and $\ell$ are their multiplicities.

Let $E_\lambda$ and $E_\kappa$ denote the parallel tangent
eigendistributions.  Consider the ambient subspaces, initially depending
on $x$,
\[
 \mathcal E(x)=E_\lambda(x)\oplus\operatorname{span}\{\lambda N(x)-x\},
 \qquad
 \mathcal F(x)=E_\kappa(x)\oplus\operatorname{span}\{\kappa N(x)-x\}.
\]
They are orthogonal and their dimensions sum to $m+2$, because
$\langle\lambda N-x,\kappa N-x\rangle=1+\lambda\kappa=0$.
For any tangent vector $w$ and any local section $V$ of $E_\lambda$,
ambient differentiation gives
\[
 D_wV=\nabla_wV+\langle w,V\rangle(\lambda N-x),
 \qquad
 D_w(\lambda N-x)=-(\lambda A+I)w.
\]
Both expressions lie in $\mathcal E(x)$: the second vanishes on
$E_\kappa$ and maps $E_\lambda$ to itself.  Hence $\mathcal E$ is a
constant subspace of $\mathbb R^{m+2}$ along the connected hypersurface;
the same holds for its orthogonal complement $\mathcal F$.
The projections of $x$ to these fixed spaces have lengths
\[
 |\pi_{\mathcal E}x|=\frac1{\sqrt{1+\lambda^2}}
                    =\sqrt{\frac{k}{m}},
 \qquad
 |\pi_{\mathcal F}x|=\frac1{\sqrt{1+\kappa^2}}
                    =\sqrt{\frac{\ell}{m}}.
\]
Thus the image of $\Sigma$ lies in the displayed Clifford product.
The inclusion is a local diffeomorphism onto this product, since both
manifolds have dimension $m$.  Its image is open, and compactness makes
it closed.  The product is connected because $k,\ell\ge1$, so the image
is the entire product.  The final assertion follows from the definition
of $a_m$.
\end{proof}
